\documentclass[11pt]{article}

\usepackage[a4paper, margin=1in]{geometry}
\usepackage{times}
\usepackage{hyperref}
\usepackage{enumitem}
\usepackage{parskip}
\usepackage{titlesec}
\usepackage{amsmath}

\hypersetup{colorlinks=true, linkcolor=blue, urlcolor=blue, citecolor=blue}

\titleformat{\section}{\normalfont\large\bfseries}{\thesection}{0.6em}{}
\titlespacing*{\section}{0pt}{12pt}{6pt}

\setlength{\parindent}{0pt}
\setlength{\parskip}{4pt}

\title{\textbf{Master's Thesis Proposal}\\[0.3em]
Visual Analytics for Atomic Fact Annotation:\\
Guideline Refinement through a Conflict-Aware Knowledge Graph}
\author{Zilong Yan \\[0.2em]
\quad First Reviewer: Prof.\ Daniel A.\ Keim}
\date{3 May 2026}

\begin{document}
\maketitle

\section{Motivation}

Factuality evaluation of Large Language Models (LLMs) increasingly relies on
decomposing generated text into discrete \textit{atomic facts}: minimal,
self-contained units of meaning against which model outputs can be verified.
Pipelines such as FActScore \cite{min2023factscore} and FEVER
\cite{thorne2018fever} treat these atomic facts as the granule of evaluation.
However, no agreed definition of atomicity exists, and recent empirical work
shows substantial inter-annotator disagreement (IAA) among both human experts
and instruction-tuned LLMs \cite{schmidt2025dissecting,wanner2024closer},
driven by ambiguity in coreference resolution, conditional decomposition,
and granularity trade-offs.
Schmidt et al.\ \cite{schmidt2025dissecting}
argue that low IAA is a signal of underspecified annotation guidelines and
propose a Visual Analytics (VA) workflow to surface disagreement and drive
iterative guideline refinement.

A practical obstacle to this refinement loop is that presenting disagreements as a flat list (e.g., a spreadsheet) is not actionable.
Without visual structure, an analyst cannot easily answer concrete questions needed to fix guidelines: Are models consistently disagreeing about dates?
Is confusion centered around specific organizations? Or does a particular legal clause always trigger conflicts?
Furthermore, surfacing genuine disagreement at scale requires comparing facts extracted by multiple annotators.
Using pairwise LLM-based comparison for this arbitration requires $O(n^2)$ LLM calls, scaling quadratically and becoming prohibitively expensive.

This thesis proposes a Visual Analytics prototype centred on a
\textit{conflict-aware knowledge graph} (KG): an analyst-facing structure in
which inter-model disagreements become visually discoverable, drawing
attention precisely to regions where the annotation guideline is
underspecified (e.g., entity boundary ambiguity or conditional clause trigger conflicts).
A two-layer conflict detection mechanism---combining a fast 
rule-based text filter with a deeper LLM arbitrator---serves the KG by
filtering redundant fact pairs locally and reserving LLM arbitration for
genuinely ambiguous cases, keeping the system responsive for
interactive experimentation.
Inside the tool, an interactive 
guideline editor closes the refinement loop: analysts revise extraction rules directly within the interface and immediately observe 
the resulting shift in the KG.
The principal evaluation is a closed-loop 
self-experiment that measures whether one round of KG-driven revision 
(\texttt{v1} $\to$ \texttt{v2}) shifts the conflict-type distribution 
toward inter-model convergence.
Minor adjustments to the corpus or model set may be made in consultation with the supervisor, while the overall experimental framework remains fixed.

\section{Goals and Methodology}

The project is organized into four distinct phases.

\subsection*{Phase 1: Data Pipeline \& Decontextualized Fact Extraction}

We adopt EUR-Lex English-language EU regulatory documents as the primary
corpus, selected for structured legal prose, public availability, and
linguistic compatibility with mid-sized open-weight models.
The experimental set is sized at 50--80 documents so that several thousand 
aligned pairs are available for distribution-level comparisons in Phase 4.

Three 4B-class local models are deployed via Ollama:
\begin{itemize}[leftmargin=1.5em, itemsep=2pt]
    \item \textbf{Qwen3.5-4B} \cite{qwen2026}: strong multilingual reasoning baseline.
    \item \textbf{Gemma3-4B} \cite{gemma2026}: stable English instruction-following.
    \item \textbf{Phi4-mini (3.8B)} \cite{phi2026}: optimized for structured output.
\end{itemize}

These recent small models were selected because technical reports 
demonstrate their strong performance in structured JSON generation and single-hop reasoning, 
achieving competitive results in constrained extraction settings \cite{qwen2026, phi2026}.
The extraction prompt enforces decontextualization following Molecular Facts \cite{gunjal2024molecular}: each fact must be free of unresolved pronouns and carry a \textit{minimal but sufficient} set of disambiguating modifiers --- role, date, location --- such that it remains uniquely identifiable without source context.
The \textit{minimum-constraint} formulation is adopted directly from \cite{gunjal2024molecular},
which also motivates treating the level of constraint as a
\textit{guideline-controlled parameter}: different guideline rules induce different
effective decontextualization regimes, and Phase 4 measures the consequences.
Each fact carries a \texttt{source\_quote} field enabling character-level offset recovery via RapidFuzz, so that every triple in the KG can be traced
back to a specific span in the source document.
Output is constrained by a \textbf{source-agnostic} JSON Schema. It mitigates formatting failures and serves as a universal interface to seamlessly ingest externally provided \textit{human-annotated} facts without requiring a custom UI.

\subsection*{Phase 2: Semantic Alignment \& Two-Layer Conflict Detection}

\textbf{Architectural Justification.} The decision to adopt a ``decouple and align'' 
pipeline is strictly guided by empirical findings from our preceding benchmarking project \cite{yan2026benchmarking}.
Extensive evaluations revealed that local 4B to 7B-class models exhibit severe 
degradation in effective context windows---showing significant performance drops 
when processing beyond 5,000 tokens---and systematically struggle with multi-hop 
reasoning tasks across aggregated documents \cite{yan2026benchmarking}.
This phenomenon aligns with broader literature documenting the long-context limitations 
of small transformers \cite{kuratov2024babilong, wang2025longcontext}.
By reducing the task to local extraction followed by pairwise atomic comparison, we constrain 
the LLM input to 1K--2K tokens per task, successfully bypassing these bottlenecks.

Facts from different models are aligned using Sentence-BERT \cite{reimers2019sbert}
combined with the Hungarian algorithm for optimal one-to-one assignment \cite{schmidt2025dissecting}.
A default similarity threshold of 0.78 is adopted initially; this will be dynamically adjusted based on downstream performance.
The unit of analysis is the \textit{aligned pair} (two extracted facts mapping to the same underlying semantic claim). This avoids the ambiguity inherent in strict entity-level KG comparison, where surface-form variations and differing graph topologies create false structural mismatches.

Aligned pairs are routed through a two-layer detection mechanism:
\begin{itemize}[leftmargin=1.5em, itemsep=2pt]
    \item \textbf{Layer 1 (rule-based filter)}: Pairs with cosine
        similarity $> 0.95$ and high object-field character
        overlap are labelled \textsc{Redundancy} and bypass LLM arbitration.
        Pairs with high sentence-level similarity but
        divergent object values (the \textit{antonym trap},
        e.g., ``rate is 15\%'' vs.\ ``rate is 20\%'') escalate to Layer 2.
    \item \textbf{Layer 2 (LLM arbitration)}: Remaining pairs are sent
        to Qwen3.5-4B with a few-shot prompt
        returning \textsc{Contradiction}, \textsc{Granularity},
        \textsc{Redundancy}, or \textsc{No\_Conflict}.
\end{itemize}

\textbf{1-Hop Limitation.} It must be explicitly noted that this mechanism is constrained to \textbf{1-hop semantic resolution}. It evaluates direct contradictions between individual aligned pairs, but cannot autonomously detect multi-hop structural conflicts spanning multiple nodes in the wider knowledge graph.

This mechanism is positioned as infrastructure to keep per-document processing fast enough for interactive guideline revision.

\subsection*{Phase 3: Conflict-Aware Knowledge Graph \& VA Interface}

This phase delivers the central contribution: a \textit{conflict-aware knowledge graph} that groups facts around shared entities and exposes inter-model disagreement as visual structure --- the form in which guideline-relevant patterns become discoverable.

\textbf{Triple Extraction, Provenance, and Graph Storage.} Each atomic fact is converted
into an (S,~P,~O) triple augmented with provenance metadata (source model, document ID, offsets, conflict label).
To enable efficient structural queries and overcome the 1-hop limitation, these triples are persisted into a \textbf{graph database (e.g., Neo4j)}.
This storage backend provides the infrastructure to quickly identify topological inconsistencies and execute complex pattern-matching algorithms.

\textbf{Entity resolution.} Surface-form variants
are unified via SBERT embedding clustering over the global subject-object pool, with
an editable merge threshold.
Without this step, the KG fragments and disagreement signals are lost.

\textbf{Three-level node colouring and edge typing.} Nodes inherit the
strongest conflict label across incident triples:
\textit{green} (all agree),
\textit{yellow} (\textsc{Granularity}/\textsc{Redundancy}),
\textit{red} (\textsc{Contradiction}).
Edges are typed by predicate and styled by conflict status.

\textbf{Visual analytics interactions.} The front-end is built with Cytoscape.js over a Python (Flask) back-end. Core interactions:
\begin{itemize}[leftmargin=1.5em, itemsep=2pt]
    \item \textit{Click-to-highlight}: Clicking a node surfaces supporting text spans.
    \item \textit{Conflict-cluster filtering}: Filter the KG to show specific disagreement patterns.
    \item \textit{In-tool fact verification}: Clicking a red node opens a pre-populated search URL.
    \item \textit{Interactive Graph Constraints ("Patching"):} To compensate for the LLM's multi-hop reasoning limitations, analysts can inject domain-specific structural rules directly via the UI (e.g., setting a cardinality constraint such as ``a specific legal sub-clause can belong to at most one parent clause''). The Neo4j backend queries the KG against these constraints. Topological violations are surfaced as \textsc{Contradiction}, allowing human analysts to effectively "patch" the system's reasoning capabilities.
    \item \textit{Interactive guideline live-edit panel}: Revise rules and re-run extraction locally.
\end{itemize}

The KG is the analytical surface through which guideline diagnosis happens.
A key limitation is the ``hen and egg'' problem of using LLMs to evaluate LLM outputs: absolute factual guarantees are impossible.
Thus, Phase 4 measures the effect of using this surface to drive refinement, not the KG itself.

\subsection*{Phase 4: Iterative Guideline Refinement Experiment}

The principal evaluation is a closed-loop self-experiment of
\textit{iterative guideline refinement}.
Additional quantitative analyses
will be conducted as concrete questions emerge from Phases 1--3.

\textbf{Iterative Guideline Refinement (closed-loop).} A baseline guideline (\texttt{v1}) processes the entire EUR-Lex corpus.
KG diagnostics will be analyzed to identify high-frequency disagreement patterns.
A targeted revision (\texttt{v2}) is then authored to address these patterns, and the pipeline is re-run.

\textbf{Scope of Evaluation:} The primary quantitative evaluation remains strictly \textbf{inter-model}.
However, to demonstrate applicability to broader inter-annotator workflows, we will ingest a small set of pre-annotated human facts via the standard JSON interface.
This serves strictly as a qualitative validation to prove the tool can surface discrepancies between human ground-truth and model extractions.

Comparison is performed at the \textit{distribution level}. Reported metrics:
\textit{Conflict Type Distribution Shift},
\textit{Total Conflict Count Reduction}, and
\textit{Layer~1 Filter Rate Change}.
Tool effectiveness is demonstrated by directional improvements toward
convergence and a quantifiable reduction in contested facts a hypothetical analyst would need to review.

\begin{thebibliography}{12}

\bibitem{min2023factscore}
S.~Min, K.~Krishna, X.~Lyu, M.~Lewis, W.~Yih, P.~W.~Koh, M.~Iyyer,
L.~Zettlemoyer, and H.~Hajishirzi.
\textit{FActScore: Fine-grained Atomic Evaluation of Factual Precision in
Long Form Text Generation.} EMNLP 2023, pp.\ 12076--12100.

\bibitem{thorne2018fever}
J.~Thorne, A.~Vlachos, C.~Christodoulopoulos, and A.~Mittal.
\textit{FEVER: a Large-scale Dataset for Fact Extraction and VERification.}
NAACL 2018, pp.\ 809--819.

\bibitem{schmidt2025dissecting}
M.~Schmidt, D.~A.~Keim, and F.~L.~Dennig.
\textit{Dissecting Atomic Facts: Visual Analytics for Improving Fact
Annotations in Language Model Evaluation.} arXiv:2509.01460, 2025.

\bibitem{wanner2024closer}
M.~Wanner, S.~Ebner, Z.~Jiang, M.~Dredze, and B.~Van Durme.
\textit{A Closer Look at Claim Decomposition.}
*SEM 2024, pp.\ 153--175.

\bibitem{gunjal2024molecular}
A.~Gunjal and G.~Durrett.
\textit{Molecular Facts: Desiderata for Decontextualization in LLM Fact
Verification.} EMNLP Findings 2024, pp.\ 3751--3768.

\bibitem{reimers2019sbert}
N.~Reimers and I.~Gurevych.
\textit{Sentence-BERT: Sentence Embeddings using Siamese BERT-Networks.}
EMNLP 2019, pp.\ 3982--3992.

\bibitem{qwen2026}
Qwen Team.
\textit{Qwen3.5 Technical Report.}
arXiv preprint, 2026.

\bibitem{gemma2026}
Google DeepMind.
\textit{Gemma 3: Open Models Based on Gemini Research and Technology.}
arXiv preprint, 2026.

\bibitem{phi2026}
Microsoft Research.
\textit{Phi-4: Small Language Models for Complex Structured Output.}
arXiv preprint, 2026.

\bibitem{yan2026benchmarking}
Z.~Yan.
\textit{Benchmarking Knowledge Graph Serializations for LLM Structure-Aware Reasoning 
and the Shift Towards Text-Level Atomic Fact Resolution.}
Master's Project Report, University of Konstanz, 2026.

\bibitem{kuratov2024babilong}
Y.~Kuratov, A.~Bulatov, P.~Anokhin, I.~Rodkin, D.~Sorokin, A.~Sorokin, and M.~Burtsev.
\textit{BABILong: Testing the Limits of LLMs with Long Context Reasoning-in-a-Haystack.}
arXiv preprint arXiv:2406.10149, 2024.

\bibitem{wang2025longcontext}
L.~Wang et al.
\textit{How Well Do LLMs Handle Long Context? A Systematic Benchmark.}
arXiv preprint arXiv:2501.14427, 2025.

\end{thebibliography}

\end{document}